\section{Generations}\label{sec:generations}

In four dimensions the effective action of a step is not turned into a single activity of the new scale. As in Ba\l aban's analysis \cite[(0.22)--(0.23)]{BalYM4I}, it is a sum of generations (Remark~\ref{rem:flow-status}). The generation of scale $j$ is created by the step from scale $j-1$ to scale $j$. It is a family of activities on the polymers of scale $j$, bounded on the domains of scale $j$ only. At every later scale $k$ it is evaluated at the background field $U_k(V)$ of that scale, and only this argument changes.

This section defines a generation and its norm (Definition~\ref{def:generation}). The definition does not refer to the flow of the couplings: the couplings enter only through the small-field threshold at which the symmetry of the generation is required (Remark~\ref{rem:generation-couplings}). Lemma~\ref{lem:later-small-field} shows that a background field of a later scale is a background field of scale $j$, of a configuration in the small-field set of scale $j$. So the symmetry of a generation, a property of its sum at scale $j$, persists at all later scales. A generation is moreover covariant, as a family, under the translations and reflections of its grid, (G4). The step does not produce covariant families. Section~\ref{subsec:gen-symmetrization} averages the family it produces over the symmetries of the grid; this changes neither its sum nor its bound, but it needs a slightly smaller domain. The coefficient $\beta_j$ of the counterterm $-\beta_jS$ that accompanies a generation in the flow is not part of the definition. It is read off from the generation through its period density, a quadratic form in a constant curvature (Section~\ref{subsec:gen-period}). By (G4), the period density of a generation is a multiple $\beta_E$ of that of the action, with no topological and no anisotropic part (Proposition~\ref{prop:period-iso}).

\paragraph{Setting.} The scales are $0\le j\le K$. The torus and the parameters are as in Sections~\ref{subsec:tf-scales} and~\ref{subsec:flow-order}, and objects of scale $j$ are read in the units of scale $j$ unless said otherwise. Then $M\mid n_j$, so the polymer grid of scale $j$ is defined (Section~\ref{subsec:flow-norm}). Moreover $n_j\ge n_K\ge n_*$, and $n_*$ contains the torus conditions of Theorem~\ref{thm:balaban}, of Proposition~\ref{prop:sharp-curvature} with $\delta=1$, and of Part~VII for the step from scale $j-1$ to scale $j$ (Section~\ref{subsec:tf-scales}).

\paragraph{Grid symmetries and the generation domains.} In the units of scale $j$, let $\mathfrak S^{\rm grid}_j$ be the group of the maps $r(x)=Rx+a$ of $\Tf$ with $R\in B_4$ and $a\in M\Z^4$. It is finite. It is generated by the translations by the vectors of $M\Z^4$ and by $B_4$ acting about the centre $0$ of a grid cube of scale $j$, and it contains $B_4$ acting about the centre $c\in M\Z^4$ of any grid cube, $x\mapsto c+R(x-c)$. It consists of the elements of $\mathfrak S_j$ (Section~\ref{sec:averaging}) that map the grid cubes of scale $j$ onto grid cubes. Indeed, the grid cube centred at $c$ is the set of the points at $\ell^\infty$-distance at most $\frac12(M-\eta)$ from $c$ (Lemma~\ref{lem:grid-nesting}(a)), the elements of $\mathfrak S_j$ are $\ell^\infty$-isometries, and for $r(x)=Rx+v$ with $v\in\T^{(j)}$ we have $rc\in M\Z^4$ for all $c\in M\Z^4$ if and only if $v\in M\Z^4$. Every $r\in\mathfrak S^{\rm grid}_j$ extends to an isometry of the continuum torus that maps the closed cubes filled by the grid cubes onto each other. So it maps polymers of scale $j$ onto polymers, and it preserves $d_j$, the number $|X|_M$ of cubes and the diameter; in particular it maps wrapping polymers onto wrapping polymers, and $(rX)^{(3)}=r(X^{(3)})$. It acts on configurations by $(rU)(b)=U(r^{-1}b)$, with $U(-b)=U(b)^{-1}$, as in Section~\ref{sec:averaging}; the translation by $a$ is the map $\tau_a$ of (G2) below.

With the domain constants $\theta_r$, $\delta_1$, $\delta_0$ of Remark~\ref{rem:domain-constants} put
\[
\delta_1^\flat=\tfrac14\delta_1,\qquad\mathcal U^\flat(X)=\mathcal U(X;\theta_r,\delta_1^\flat,\delta_0)\subset\mathcal U(X)=\mathcal U(X;\theta_r,\delta_1,\delta_0).
\]
Generations are bounded on the domains $\mathcal U^\flat(X)$. The step produces families bounded on the domains $\mathcal U(X)$, and their averages over $\mathfrak S^{\rm grid}_j$ are bounded on $\mathcal U^\flat(X)$ (Section~\ref{subsec:gen-symmetrization}). The reflections in $\mathfrak S^{\rm grid}_j$ need not map $\mathcal U(X)$ into $\mathcal U(rX)$, since Definition~\ref{def:graded} uses forward differences and coordinate-ordered paths, but they map $\mathcal U^\flat(X)$ into $\mathcal U(rX)$ under a further condition (DC6) on the domain constants (Lemma~\ref{lem:grid-domains} and \eqref{eq:flat-in-full}).

\paragraph{Background fields at an admissible threshold.} Let $h$ be an admissible threshold (Section~\ref{subsec:flow-parameters}), let $j\le K$, and let $V\in\mathfrak U_j(h)$ be real. Since $h\le a_1$, the configuration $V$ has a minimal orbit (Theorem~\ref{thm:balaban}), a $\mathfrak G_0$-orbit of real configurations on $\Tf$; we write $U_j(V)$ for any of its elements. By Proposition~\ref{prop:sharp-curvature} with $\delta=1$ and Lemma~\ref{lem:critical-current}, $U_j(V)$ has curvature at most $(C_*+1)h$ and current at most $C_J(C_*+1)h$, both less than $\theta_r$. So
\[
U_j(V)\in\mathcal U^\flat(X)\subset\mathcal U(X)\qquad\text{for every polymer $X$ of scale $j$},
\]
with real part $U_j(V)$ and $A=0$ in Definition~\ref{def:graded}; the field $A=0$ satisfies the bounds of every graded domain. The set $\mathfrak U_j(h)$ is invariant under real gauge transformations and under $\mathfrak S_j$ (Section~\ref{sec:averaging}). Let $v$ be a real gauge transformation on $\T^{(j)}$, and $u$ a real one on $\Tf$ with $u|_{\T^{(j)}}=v$. Then $U\mapsto U^u$ maps the minimal orbit of $V$ onto that of $V^v$. Indeed, it maps $\mathfrak B(V)\cap\mathfrak U(\theta_0)$ onto $\mathfrak B(V^v)\cap\mathfrak U(\theta_0)$, preserves $S$ and maps critical orbits onto critical orbits (the paragraph after Theorem~\ref{thm:variational}), and the critical orbit is unique (Theorem~\ref{thm:balaban}(a)).

\begin{definition}[Generation]\label{def:generation}
Let $1\le j\le K$, and let $h$ be an admissible threshold. A \emph{generation of scale $j$ at threshold $h$} is a family $E=\{E(X,\cdot)\}$, indexed by the polymers $X$ of scale $j$ (Section~\ref{subsec:flow-norm}), with the following properties, in the units of scale $j$.
\begin{enumerate}[label=(G\arabic*)]
\item \emph{(Activities.)} Each $E(X,\cdot)$ is a normalized graded activity at $X$ on $\mathcal U^\flat(X)=\mathcal U(X;\theta_r,\delta_1^\flat,\delta_0)$ (Definition~\ref{def:graded}), with the domain constants of Remark~\ref{rem:domain-constants} and $\delta_1^\flat=\frac14\delta_1$, and $E(X,U)$ is real for real $U\in\mathcal U^\flat(X)$. So $E(X,U)$ depends only on $U$ on the bonds of $X$, it is holomorphic and invariant under real gauge transformations on $\mathcal U^\flat(X)$, and $E(X,1)=0$.
\item \emph{(Grid covariance.)} For every $a\in M\Z^4$, every polymer $X$ and every $U\in\mathcal U^\flat(X)$,
\[
E(X+a,\tau_aU)=E(X,U),\qquad(\tau_aU)(b)=U(b-a).
\]
\item \emph{(Invariance of the sum.)} The function
\[
\Phi_E(V)=\sum_XE\big(X,U_j(V)\big),\qquad V\in\mathfrak U_j(h),
\]
satisfies $\Phi_E(rV)=\Phi_E(V)$ for every $V\in\mathfrak U_j(h)$ and every $r\in\mathfrak S_j$.
\item \emph{(Covariance under the grid symmetries.)} For every $r\in\mathfrak S^{\rm grid}_j$, every polymer $X$ and every $U\in\mathcal U^\flat(X)$ with $rU\in\mathcal U^\flat(rX)$,
\[
E(rX,rU)=E(X,U),\qquad(rU)(b)=U(r^{-1}b).
\]
\end{enumerate}
The \emph{norm} of $E$ is
\[
\|E\|_j=\sup_Xe^{\kappa_nd_j(X)}\|E(X)\|,\qquad\|E(X)\|=\sup_{\mathcal U^\flat(X)}|E(X,\cdot)|,
\]
with $\kappa_n$ of Remark~\ref{rem:parameters}(1) and $d_j$ the tree length of the grid of scale $j$ (Section~\ref{subsec:flow-norm}).
\end{definition}

In (G2), $\tau_a$ maps $\mathcal U^\flat(X)$ onto $\mathcal U^\flat(X+a)$, since translations preserve curvatures, currents, covariant differences and the coordinate-ordered paths of Definition~\ref{def:graded}. So (G2) is the case of the translations in (G4); we keep it separately, since the cell sums of the period density use only (G2) (Lemma~\ref{lem:period-density}(c)). For the other elements of $\mathfrak S^{\rm grid}_j$ the condition $rU\in\mathcal U^\flat(rX)$ in (G4) is not automatic (Lemma~\ref{lem:grid-domains}). In (G3) the sum is finite, and $\Phi_E$ is well defined by the paragraph before the definition. Each $U_j(V)$ lies in every $\mathcal U^\flat(X)$, and $\Phi_E(V)$ does not depend on the element $U_j(V)$ of the minimal orbit, since the orbit is a $\mathfrak G_0$-orbit and each $E(X,\cdot)$ is invariant under real gauge transformations. For the same reason $\Phi_E(V^v)=\Phi_E(V)$ for real gauge transformations $v$ on $\T^{(j)}$. By (G1), $\Phi_E$ is real. The domains $\mathcal U^\flat(X)$ are smaller than the domains $\mathcal U(X)$ of the step because of the symmetrization (Section~\ref{subsec:gen-symmetrization}).

\begin{remark}[The couplings in the definition]\label{rem:generation-couplings}
The couplings enter Definition~\ref{def:generation} only through the threshold $h$, and only in (G3). In the flow, the generation of scale $j$ is taken at the threshold $h_j=p(g_j)g_j$ of the density of scale $j$ (Section~\ref{subsec:flow-parameters}). Nothing in the definition depends on how the couplings $g_0,\dots,g_j$ are chosen. They may be the couplings of the flow, defined by the recursion $g_{i+1}^{-2}=g_i^{-2}-\beta_{i+1}$, or an arbitrary sequence in $(0,\gamma]$, as for the continuity in $g_0,\dots,g_j$ (Proposition~\ref{prop:flow-continuity}) required in (F1) of Section~\ref{sec:ol-uv}. The step that produces the generations needs more: (Q2) requires $g_{k+1}\le2g_k$ (Remark~\ref{rem:flow-constants}, consequences), as does Lemma~\ref{lem:later-small-field}(c). Continuity in the couplings is therefore to be stated on sequences with $g_{i+1}\le2g_i$. Since $\mathfrak U_j(h')\subset\mathfrak U_j(h)$ for $h'\le h$, and both sets are invariant under $\mathfrak S_j$, a generation at threshold $h$ is one at every threshold $h'\le h$. The norm does not depend on the threshold.
\end{remark}

The following lemma is the input for (G3) at later scales. Part~(c) is the only place where the couplings enter.

\begin{lemma}[Background fields at later scales]\label{lem:later-small-field}
Let $1\le j\le k\le K$, let $h$ be an admissible threshold, and let $V\in\mathfrak U_k(h)$ be a real configuration on $\T^{(k)}$, in the units of scale $k$. Let $U=U_k(V)$ be an element of its minimal orbit, a configuration on $\Tf$, and, for $j\le i\le k$, let $U^{(i)}$ be its average on $\T^{(i)}$ (Definition~\ref{def:one-step}). Then, for $j\le i\le k$:
\begin{enumerate}[label=(\alph*)]
\item In the units of scale $i$, $U^{(i)}\in\mathfrak U_i\big(2(C_*+1)L^{-2(k-i)}h\big)$ if $i<k$, and $U^{(k)}=V$. In particular $U^{(i)}\in\mathfrak U_i(h)$.
\item $U$ is an element of the minimal orbit of $U^{(i)}$ for the problem of scale $i$; that is, one may take $U_i(U^{(i)})=U$.
\item Let $g_j,\dots,g_k\in(0,\gamma]$ with $g_{i+1}\le2g_i$ for $j\le i<k$, and let $h=h_k$. Then $2(C_*+1)L^{-2(k-j)}h_k\le h_j$ if $j<k$, and $U^{(j)}\in\mathfrak U_j(h_j)$.
\item Let $E$ be a generation of scale $j$ at an admissible threshold $h'$, with $h'\ge h$ if $j=k$, and $h'\ge2(C_*+1)L^{-2(k-j)}h$ if $j<k$; under the hypotheses of~(c), $h'=h_j$ is such a threshold. Then $U\in\mathcal U^\flat(X)$ for every polymer $X$ of scale $j$, in the units of scale $j$, and
\[
\sum_XE\big(X,U_k(V)\big)=\Phi_E\big(U_k(V)^{(j)}\big).
\]
The left side does not depend on the element $U_k(V)$ of the minimal orbit. As a function of $V\in\mathfrak U_k(h)$ it is invariant under real gauge transformations and under $\mathfrak S_k$.
\end{enumerate}
\end{lemma}

\begin{proof}\hasdetail{lem:later-small-field}
By the Setting, the torus conditions of Theorem~\ref{thm:balaban} and Proposition~\ref{prop:sharp-curvature} hold at every scale $i\le K$, and those of Part~VII hold for every step from scale $i$ to scale $i+1\le K$.

(a) Since $U\in\mathfrak B(V)$, $U^{(k)}=V$. Let $i<k$. By Proposition~\ref{prop:sharp-curvature} with $\delta=1$, which applies since $h\le a_{\delta=1}$, $U\in\mathfrak U_0((C_*+1)h)$ in the units of scale $k$. Since $(C_*+1)h\le c_1$, Proposition~\ref{prop:avg-regular} gives $U^{(i)}\in\mathfrak U_i(2(C_*+1)h)$, in the same units. The curvature of a configuration $Y$ on $\T^{(i)}$ is $\delta_i^{-2}\ell(Y(\partial p))$ (Definition~\ref{def:curvature}), where the spacing $\delta_i$ of $\T^{(i)}$ is $L^{i-k}$ in the units of scale $k$ and $1$ in those of scale $i$. So in the units of scale $i$ the bound is multiplied by $L^{-2(k-i)}$. Finally $2(C_*+1)L^{-2(k-i)}\le2(C_*+1)L^{-2}\le\frac14$, by $L^2\ge8(C_*+1)$ (Remark~\ref{rem:parameters}(2)).

(b) By downward induction on $i$. For $i=k$ this is the choice of $U$. Let $j\le i<k$, and suppose that $U$ is an element of the minimal orbit of $U^{(i+1)}$ for the problem of scale $i+1$. Apply Lemma~\ref{lem:composition} to the step from scale $i$ to scale $i+1$, in the units of scale $i$ (so $\T=\T^{(i)}$ and $\T'=\T^{(i+1)}$), with $V'=U^{(i+1)}$, $\eps'_1=h$ and $U'=U$. Its hypotheses hold: $V'$ is real, its curvature in the units of scale $i+1$ is at most $h$ by (a), and $h$ satisfies the conditions on $\eps'_1$, since it is admissible. The base configuration of the lemma is $V^0=U'^{(i)}=U^{(i)}$, and part~(b) of the lemma states that $U$ is an element of the minimal orbit of $U^{(i)}$ for the problem of scale $i$.

(c) By the paragraph on consecutive thresholds in Section~\ref{subsec:flow-parameters}, $h_{i+1}\le2h_i$ for $j\le i<k$, so $h_k\le2^mh_j$ with $m=k-j$. The threshold $h_k$ is admissible, since $g_k\in(0,\gamma]$. If $m=0$, then $U^{(j)}=V\in\mathfrak U_j(h_j)$. If $m\ge1$, then, since $2L^{-2}\le1$ and $L^2\ge8(C_*+1)$,
\[
2(C_*+1)L^{-2m}h_k\le2(C_*+1)\big(2L^{-2}\big)^mh_j\le4(C_*+1)L^{-2}h_j\le\tfrac12h_j ,
\]
and (a) gives $U^{(j)}\in\mathfrak U_j(\frac12h_j)\subset\mathfrak U_j(h_j)$.

(d) By (a) and the hypothesis on $h'$, $U^{(j)}\in\mathfrak U_j(h')$. Under the hypotheses of (c), $h_j$ is an admissible threshold, since $g_j\in(0,\gamma]$, and it satisfies the hypothesis on $h'$, by the bound of (c). By (b), $U$ is an element of the minimal orbit of $U^{(j)}$. Since $h'$ is admissible, the paragraph before Definition~\ref{def:generation} gives $U\in\mathcal U^\flat(X)$ for every polymer $X$ of scale $j$, and the definition of $\Phi_E$, with $U_j(U^{(j)})=U$, gives the identity. Two elements of the minimal orbit of $V$ differ by a real gauge transformation, and each $E(X,\cdot)$ is invariant under real gauge transformations on $\mathcal U^\flat(X)$, which is invariant under them. So the left side does not depend on the element. Let $v$ be a real gauge transformation on $\T^{(k)}$, and $u$ a real one on $\Tf$ with $u|_{\T^{(k)}}=v$. By the paragraph before Definition~\ref{def:generation}, at scale $k$, $V^v\in\mathfrak U_k(h)$ and $U^u$ is an element of the minimal orbit of $V^v$; and $E(X,U^u)=E(X,U)$ for every $X$. This proves the gauge invariance. Let $r\in\mathfrak S_k$. By Theorem~\ref{thm:balaban}(c), $rV\in\mathfrak U_k(h)$, and $rU$ is an element of the minimal orbit of $rV$. Since $\mathfrak S_k\subset\mathfrak S_j\subset\mathfrak S_i$ for $i\le j$, Lemma~\ref{lem:one-step-basic}(d), applied to the averagings $\mathfrak a_1,\dots,\mathfrak a_j$ in the units of scale $k$, gives $(rU)^{(j)}=rU^{(j)}$. The identity for $rV$, then (G3) for $U^{(j)}\in\mathfrak U_j(h')$ and $r\in\mathfrak S_j$, then the identity for $V$ give
\[
\sum_XE(X,rU)=\Phi_E\big(rU^{(j)}\big)=\Phi_E\big(U^{(j)}\big)=\sum_XE(X,U) .\qedhere
\]
\end{proof}

\begin{remark}[Symmetry of the density at later scales]\label{rem:later-symmetry}
Let $g_1,\dots,g_k\in(0,\gamma]$ with $g_{i+1}\le2g_i$ for $i<k$, and let $\rho_k$ have the form of Remark~\ref{rem:flow-status} on $\mathfrak U_k(h_k)$, and vanish outside it, with generations $E^{(j)}$ of scale $j$ at the thresholds $h_j$, $1\le j\le k$. Then $\rho_k(rV)=\rho_k(V)$ for every configuration $V$ and every $r\in\mathfrak S_k$. Indeed, $\mathfrak U_k(h_k)$ is invariant under $\mathfrak S_k$, the terms $S(U_k(V))$ are invariant by Theorem~\ref{thm:balaban}(c), and the sums of the generations by Lemma~\ref{lem:later-small-field}(c),(d). Since $\mathfrak S_{k+1}\subset\mathfrak S_k$, this is the hypothesis $\rho\circ r=\rho$ of the symmetry statement of Theorem~\ref{thm:sf-step}, for every $r\in\mathfrak S_{k+1}$.
\end{remark}

\subsection{Symmetrization}\label{subsec:gen-symmetrization}

The step produces families of activities whose sum has the symmetry (G3), but not the covariance (G4): single activities need not be covariant (Theorem~\ref{thm:sf-step}). This subsection averages such a family over $\mathfrak S^{\rm grid}_j$. Lemmas~\ref{lem:collar-pairs} and~\ref{lem:grid-domains} compare the domain of a polymer with that of its image under an element of $\mathfrak S^{\rm grid}_j$. Definition~\ref{def:symmetrization} and Lemma~\ref{lem:symmetrization} give the average. Throughout, $1\le j\le K$, and we use the units of scale $j$. By Lemma~\ref{lem:torus}(a), $n_j\ge8LM$.

\paragraph{Why the next two lemmas are needed.} The symmetrization evaluates the activity of a polymer $X$ at the image of a configuration under a symmetry $r$ of the grid. The graded domains of Definition~\ref{def:graded} are not invariant under these symmetries. They use forward differences, which a reflection turns into backward ones, and H\"older quotients on the collar $X^{(3)}$, whose pairs of points a symmetry need not map to pairs of the same kind near the outer boundary (Remark~\ref{rem:symmetrization-flow}(d)). Lemma~\ref{lem:grid-domains} measures the loss: the gradient budget at most doubles, up to a term $C_{\rm refl}\theta_r\delta'_0$, and the factor $2$ comes from Lemma~\ref{lem:collar-pairs}(c). Under (DC6) this gives \eqref{eq:flat-in-full}: a symmetry of the grid maps the domain $\mathcal U^\flat(X)$, with $\delta_1^\flat=\frac14\delta_1$, into the domain $\mathcal U(rX)$, on which the raw output of a step is bounded. This is why the generations are bounded on the domains $\mathcal U^\flat$ (the paragraph on (DC6) below).

\begin{lemma}[Pairs in the collar]\label{lem:collar-pairs}
Let $X$ be a polymer of scale $j$, and let $\Pi\subset\{0,\eta,2\eta\}^4$ be a set of vectors with $0\in\Pi$. Let $y\ne y'$ be points of $\Tf$ with $|y-y'|_\infty\le6$ such that $y+\Pi$ and $y'+\Pi$ lie in $X^{(3)}$. Then there is a point $q$ of $\Tf$ with the following properties.
\begin{enumerate}[label=(\alph*)]
\item For every $\mu$, $q_\mu$ lies on the shorter arc of the coordinate circle between $y_\mu$ and $y'_\mu$. Hence $|y-q|\le|y-y'|$ and $|q-y'|\le|y-y'|$, for the Euclidean, the $\ell^1$- and the $\ell^\infty$-distance.
\item $q+\Pi\subset X^{(3)}$.
\item There are points $w,w'$ of $X$ with $y,q\in\Lam_3(w)$ and $q,y'\in\Lam_3(w')$. The cubes $\Lam_3(w)$ and $\Lam_3(w')$ lie in $X^{(3)}$.
\end{enumerate}
\end{lemma}

Part~(c) is the point. The pairs $(y,q)$ and $(q,y')$ enter the H\"older quotients of Definition~\ref{def:graded} on $X^{(3)}$, while the pair $(y,y')$ need not: near the outer boundary of $X^{(3)}$, for instance near a corner at which two cubes of $X$ meet, $y$ and $y'$ need not lie in a common cube $\Lam_3(w)\subset X^{(3)}$.

\begin{proof}\hasdetail{lem:collar-pairs}
Let $\Gamma_X\subset M\Z^4$ be the set of the centres of the grid cubes of $X$, and put $R=\frac12(M-\eta)+3$. The grid cube centred at $c$ is the set of the points at $\ell^\infty$-distance at most $\frac12(M-\eta)$ from $c$ (Lemma~\ref{lem:grid-nesting}(a)). So a point $x$ of $\Tf$ lies in $X^{(3)}$ if and only if $|x-c|_\infty\le R$ for some $c\in\Gamma_X$. For a point $t$ of a coordinate circle, let $K(t)$ be the set of the points $k\in M\Z$ of the circle at distance at most $R$ from $t$. Then
\begin{equation}\label{eq:collar-membership}
x\in X^{(3)}\iff\Gamma_X\cap\big(K(x_1)\times K(x_2)\times K(x_3)\times K(x_4)\big)\ne\emptyset .
\end{equation}
The arcs $\{t:|t-k|\le R\}$, $k\in M\Z$, have length $2R=M+6-\eta<2M$, since $M\ge64$ (Remark~\ref{rem:parameters}(4)), and there are $n_j/M\ge8L$ of them on the circle. So every point lies in one or in two of them. The points that lie in two of them form the \emph{bands} $[k+M-R,k+R]$, $k\in M\Z$; a band is an arc of length $2R-M=6-\eta$, and two bands are at distance $2M-2R=M+\eta-6$. Off the bands $K$ is constant along any arc that does not meet a band, and on a band it is the union of its values just before and just after the band. The end points of the bands are points of the fine lattice, since $\frac12(M-\eta)$ and $3$ are multiples of $\eta$.

\emph{One coordinate.} Let $a,a'$ be points of the circle at distance at most $6$, let $[s,l]$ be the shorter arc between them, and let $J=[s,l+2\eta]$. We claim that there is a point $b\in[s,l]$ of the fine lattice with
\begin{equation}\label{eq:collar-choice}
K(b+\pi)\supset K(a+\pi)\cup K(a'+\pi)\qquad\text{for }\pi\in\{0,\eta,2\eta\}.
\end{equation}
Since $a+\pi$ and $a'+\pi$ lie in $[s+\pi,l+\pi]\subset J$, it suffices that $K(b+\pi)\supset K(t)$ for all $t\in[s+\pi,l+\pi]$. The arc $J$ has length at most $6+2\eta<M+\eta-6$, since $M\ge64$ (Remark~\ref{rem:parameters}(4)), so it meets at most one band. If it meets none, or if it lies in a band, $K$ is constant on $J$, and we take $b=s$. Otherwise let $B=[\beta,\beta']$ be the band that $J$ meets, so $\beta'=\beta+6-\eta$, and at least one end point of $B$ lies in $J$. If $\beta\in J$ and $\beta'\notin J$, then $K$ increases (for the inclusion) along $J$, and we take $b=l$. If $\beta'\in J$ and $\beta\notin J$, then $K$ decreases along $J$, and we take $b=s$. If $\beta,\beta'\in J$, we take $b=\beta$. Then $s\le\beta\le\beta'-2\eta\le l$, since $\beta'\le l+2\eta$ and $6-\eta\ge2\eta$. Moreover $b+\pi\in B$ for $\pi\le2\eta$, so $K(b+\pi)$ is the largest value of $K$ on $J$. In all cases \eqref{eq:collar-choice} holds.

\emph{The point $q$.} Let $q_\mu$ be the point $b$ of the first step for $a=y_\mu$ and $a'=y'_\mu$. This gives (a). Let $\pi\in\Pi$. By \eqref{eq:collar-choice}, $K(q_\mu+\pi_\mu)\supset K(y_\mu+\pi_\mu)$ for every $\mu$, and $y+\pi\in X^{(3)}$; so \eqref{eq:collar-membership} gives $q+\pi\in X^{(3)}$, which is (b). For (c), choose $c\in\Gamma_X$ with $c_\mu\in K(y_\mu)$ for all $\mu$; it exists, since $y\in X^{(3)}$. By \eqref{eq:collar-choice} with $\pi=0$, also $c_\mu\in K(q_\mu)$, so $|y-c|_\infty\le R$ and $|q-c|_\infty\le R$. In each coordinate the arcs $[\max(y_\mu,q_\mu)-3,\min(y_\mu,q_\mu)+3]$ and $[c_\mu-\frac12(M-\eta),c_\mu+\frac12(M-\eta)]$ then meet, since $|y_\mu-q_\mu|\le6$, and they meet in a point of the fine lattice. So there is a point $w$ of the grid cube centred at $c$, hence of $X$, with $|w-y|_\infty\le3$ and $|w-q|_\infty\le3$. The same argument for $q$ and $y'$, with $c'\in\Gamma_X$ chosen for $y'$, gives $w'$. Since $w,w'\in X$, the cubes $\Lam_3(w)$ and $\Lam_3(w')$ lie in $X^{(3)}$.
\end{proof}

\begin{lemma}[Grid symmetries and graded domains]\label{lem:grid-domains}
There is a constant $C_{\rm refl}\ge1$, depending only on $d$, with the following property. Let $r\in\mathfrak S^{\rm grid}_j$, let $X$ be a polymer of scale $j$, and let $\theta_r\le1$ and $0<\delta'_1\le\delta'_0\le1$. Let $U=e^{\ii\eta A}U_r$ on the bonds of $X^{(3)}$ be a representation as in Definition~\ref{def:graded} for $\mathcal U(X;\theta_r,\delta'_1,\delta'_0)$, and let $rA$ be the field on the bonds of $(rX)^{(3)}$ given by $(rA)(b)=A(r^{-1}b)$, with $A(-b)=-R(U_r(b))^{-1}A(b)$ on reversed bonds. Then $rU=e^{\ii\eta\,rA}\,rU_r$ on the bonds of $(rX)^{(3)}$, the configuration $rU_r$ is real with $|F_{rU_r}|<\theta_r$ and $|J_{rU_r}|<\theta_r$ everywhere, $|rA|=|A|$, and, with $\nabla$ taken with respect to $U_r$ and to $rU_r$,
\begin{equation}\label{eq:grid-domains}
|\nabla(rA)|_{(rX)^{(3)}}+[\nabla(rA)]_{\alpha,(rX)^{(3)}}\le|\nabla A|_{X^{(3)}}+2[\nabla A]_{\alpha,X^{(3)}}+\tfrac12C_{\rm refl}\,\theta_r\big(|\nabla A|_{X^{(3)}}+|A|\big).
\end{equation}
Here $|A|$ is the maximum over the bonds of $X^{(3)}$. Consequently
\[
r\,\mathcal U(X;\theta_r,\delta'_1,\delta'_0)\subset\mathcal U\big(rX;\theta_r,2\delta'_1+C_{\rm refl}\theta_r\delta'_0,\delta'_0\big).
\]
\end{lemma}

\begin{proof}\hasdetail{lem:grid-domains}
Write $r(x)=Rx+a$, with $Re_\mu=s_\mu e_{\pi(\mu)}$ for a permutation $\pi$ and signs $s_\mu$. For each $\mu$ put $\mu'=\pi^{-1}(\mu)$ and $\sigma_\mu=s_{\mu'}$, so that $R^{-1}e_\mu=\sigma_\mu e_{\mu'}$. We use $\|R(e^{\ii Y})-1\|\le2e^2|Y|$ for $Y\in\mathfrak g$ with $|Y|\le1$, which follows from $R(e^{\ii Y})=e^{\ii\operatorname{ad}Y}$ and $|[Y,Z]|\le2|Y||Z|$ (Section~\ref{sec:background}). In this proof $C$ denotes constants that depend only on $d$.

\emph{The real part and the sup norm.} The configuration $rU_r$ is real, and $|F_{rU_r}(p)|=|F_{U_r}(r^{-1}p)|$ and $|J_{rU_r}|_\infty=|J_{U_r}|_\infty$, as in the proof of Theorem~\ref{thm:variational}(h). With the convention on reversed bonds, $r(e^{\ii\eta A}U_r)=e^{\ii\eta\,rA}\,rU_r$ bondwise, as there. Since $R(U_r(b))$ is unitary, $|rA|=|A|$. The map $r$ maps the bonds of $X^{(3)}$ onto those of $(rX)^{(3)}$.

\emph{Differences.} Fix $\mu,\nu$ and a point $x$ such that the difference $D_\nu(rA)_\mu(x)$, taken with respect to $rU_r$, involves only points of $(rX)^{(3)}$. These points are $x+\eta(\epsilon_\mu e_\mu+\epsilon_\nu e_\nu)$, $\epsilon\in\{0,1\}^2$, if $\mu\ne\nu$, and $x$, $x+\eta e_\mu$, $x+2\eta e_\mu$ if $\mu=\nu$. Put $z=r^{-1}x$, $\kappa=\eta\one_{\{\sigma_\mu=-1\}}e_{\mu'}+\eta\one_{\{\sigma_\nu=-1\}}e_{\nu'}$ and $y=z-\kappa$. Then $r^{-1}$ maps the points involved in $D_\nu(rA)_\mu(x)$ onto those involved in $D_{\nu'}A_{\mu'}(y)$; so the latter lie in $X^{(3)}$, and $z$ is one of them. With $(D_\nu\lambda)(x)=\eta^{-1}\big(R(U(x,x+\eta e_\nu))\lambda(x+\eta e_\nu)-\lambda(x)\big)$ (Section~\ref{sec:background}), the definitions of $rA$ and $rU_r$ give
\begin{equation}\label{eq:refl-diff}
D_\nu(rA)_\mu(x)=\sigma_\mu\sigma_\nu\,T_x\,D_{\nu'}A_{\mu'}(y)+e_{\mu\nu}(x).
\end{equation}
Here $T_x=R(U_r(\gamma_x))^{-1}$ is the transport along a path $\gamma_x$ of at most two bonds from $y$ to $z$ through the points involved, and $e_{\mu\nu}(x)=0$ unless $\sigma_\mu=-1$ and $\mu\ne\nu$. We give the case $\sigma_\mu=-1$, $\sigma_\nu=1$; the case $\sigma_\mu=\sigma_\nu=-1$ is the same, with $y=z-\eta e_{\mu'}-\eta e_{\nu'}$, and in the other cases the computation gives \eqref{eq:refl-diff} with $e_{\mu\nu}(x)=0$. Let $\sigma_\mu=-1$ and $\sigma_\nu=1$, so $y=z-\eta e_{\mu'}$, and put $\hat y=y+\eta e_{\nu'}$ and $\hat z=z+\eta e_{\nu'}$. Then $(rA)_\mu(x)=-R(U_r(y,z))^{-1}A_{\mu'}(y)$, $(rA)_\mu(x+\eta e_\nu)=-R(U_r(\hat y,\hat z))^{-1}A_{\mu'}(\hat y)$ and $(rU_r)(x,x+\eta e_\nu)=U_r(z,\hat z)$. So
\[
\begin{aligned}
D_\nu(rA)_\mu(x)&=-R(U_r(y,z))^{-1}D_{\nu'}A_{\mu'}(y)\\
&\quad-\eta^{-1}\Big(R(U_r(z,\hat z))R(U_r(\hat y,\hat z))^{-1}-R(U_r(y,z))^{-1}R(U_r(y,\hat y))\Big)A_{\mu'}(\hat y).
\end{aligned}
\]
The two products in the last term are the transports from $\hat y$ to $z$ along the two pairs of sides of the plaquette $p$ with corners $y,z,\hat y,\hat z$. They differ by a factor $R(H)$ with $H$ conjugate to $U_r(\partial p)^{\pm1}=e^{\pm\ii\eta^2F_{U_r}(p)}$, and $|\eta^2F_{U_r}(p)|<\eta^2\theta_r\le1$, since $|F_{U_r}|<\theta_r$ (Definition~\ref{def:graded}). Hence
\[
|e_{\mu\nu}(x)|\le\eta^{-1}\|R(U_r(\partial p))-1\|\,|A|\le2e^2\eta\,\theta_r|A|\le2e^2\theta_r|A| .
\]
As $x$ runs over the points at which $D_\nu(rA)_\mu$ involves only points of $(rX)^{(3)}$, $y$ runs over the points at which $D_{\nu'}A_{\mu'}$ involves only points of $X^{(3)}$. So \eqref{eq:refl-diff} gives $|\nabla(rA)|_{(rX)^{(3)}}\le|\nabla A|_{X^{(3)}}+2e^2\theta_r|A|$.

\emph{Closed paths.} Let $\gamma$ be a closed lattice path of $N$ bonds in $\Tf$ with $N\eta<n_j$. Then $\|R(U_r(\gamma))-1\|\le e^2N^2\eta^2\theta_r$. Indeed, $\gamma$ lifts to a closed path in $\eta\Z^4$. Exchanging two consecutive bonds of $\gamma$ along different axes for the other two sides of the plaquette they span changes $R(U_r(\gamma))$ by at most $\|R(U_r(\partial p))-1\|\le2e^2\eta^2\theta_r$, since all transports are unitary. At most $\frac12N^2$ such exchanges sort the bonds of $\gamma$ by their axes. Since $\gamma$ is closed, its displacement along each axis vanishes; so the sorted path consists of closed paths along the coordinate lines through its initial point. Their transports are $1$, since such a path is reduced to the empty path by removing backtracks.

\emph{H\"older quotients.} Let $x\ne\tilde x$ be points of a cube $\Lam_3(\bar z)\subset(rX)^{(3)}$, and let $\mu,\nu$ be such that $D_\nu(rA)_\mu$ involves only points of $(rX)^{(3)}$ at $x$ and at $\tilde x$. Then $|x-\tilde x|_\infty\le6$, and $\eta\le|x-\tilde x|\le12$. Let $z,y$ and $\tilde z,\tilde y$ be the points defined above for $x$ and for $\tilde x$. The vector $\kappa$ is the same for both, so $\tilde y-y=R^{-1}(\tilde x-x)$ and $|\tilde y-y|=|\tilde x-x|$. Let $\Pi$ be the set of the vectors from $y$ to the points involved in $D_{\nu'}A_{\mu'}(y)$. Then $\Pi\subset\{0,\eta,2\eta\}^4$, $0\in\Pi$, and $y+\Pi$, $\tilde y+\Pi$ lie in $X^{(3)}$. Let $q,w,w'$ be the points of Lemma~\ref{lem:collar-pairs} for $y$ and $\tilde y$. Write $D(\cdot)=D_{\nu'}A_{\mu'}(\cdot)$ and $\Gamma=\Gamma_{x,\tilde x}$, and let $\gamma$ be the path from $y$ to $\tilde y$ that runs along $\gamma_x$ to $z$, along $r^{-1}\Gamma$ to $\tilde z$, and back along $\gamma_{\tilde x}$ to $\tilde y$. Since $(rU_r)(\Gamma)=U_r(r^{-1}\Gamma)$, \eqref{eq:refl-diff} gives
\[
\begin{aligned}
D_\nu(rA)_\mu(x)-R\big((rU_r)(\Gamma)\big)D_\nu(rA)_\mu(\tilde x)&=\sigma_\mu\sigma_\nu\,T_x\big(D(y)-R(U_r(\gamma))D(\tilde y)\big)\\
&\quad+e_{\mu\nu}(x)-R\big((rU_r)(\Gamma)\big)e_{\mu\nu}(\tilde x),
\end{aligned}
\]
and
\[
\begin{aligned}
D(y)-R(U_r(\gamma))D(\tilde y)&=\big(D(y)-R(U_r(\Gamma_{y,q}))D(q)\big)+R(U_r(\Gamma_{y,q}))\big(D(q)-R(U_r(\Gamma_{q,\tilde y}))D(\tilde y)\big)\\
&\quad+\big(R(U_r(\Gamma_{y,q}))R(U_r(\Gamma_{q,\tilde y}))-R(U_r(\gamma))\big)D(\tilde y).
\end{aligned}
\]
By Lemma~\ref{lem:collar-pairs}(b), $D(q)$ involves only points of $X^{(3)}$, and by (c) the pairs $(y,q)$ and $(q,\tilde y)$ lie in cubes $\Lam_3(w),\Lam_3(w')\subset X^{(3)}$. So the first two terms are at most $[\nabla A]_{\alpha,X^{(3)}}\big(|y-q|^\alpha+|q-\tilde y|^\alpha\big)\le2[\nabla A]_{\alpha,X^{(3)}}|x-\tilde x|^\alpha$; a term with $q=y$ or $q=\tilde y$ vanishes. By Lemma~\ref{lem:collar-pairs}(a), the paths $\Gamma_{y,q}$, $\Gamma_{q,\tilde y}$ and $r^{-1}\Gamma$ have length at most $|x-\tilde x|_1\le4|x-\tilde x|$, and $\gamma_x$, $\gamma_{\tilde x}$ have length at most $2\eta$. So the closed path $\Gamma_{y,q}\Gamma_{q,\tilde y}\gamma^{-1}$ has $N$ bonds with $N\eta\le12|x-\tilde x|+4\eta\le16|x-\tilde x|<n_j$, and the bound for closed paths makes the last term at most $256e^2\theta_r|x-\tilde x|^2|\nabla A|_{X^{(3)}}\le256\cdot144\,e^2\theta_r|x-\tilde x|^\alpha|\nabla A|_{X^{(3)}}$. Finally $|e_{\mu\nu}(x)|+|e_{\mu\nu}(\tilde x)|\le4e^2\eta\theta_r|A|\le4e^2\theta_r|A||x-\tilde x|^\alpha$, since $\eta\le\eta^\alpha\le|x-\tilde x|^\alpha$. Dividing by $|x-\tilde x|^\alpha$ and taking the supremum over the pairs and the components,
\[
[\nabla(rA)]_{\alpha,(rX)^{(3)}}\le2[\nabla A]_{\alpha,X^{(3)}}+C\theta_r\big(|\nabla A|_{X^{(3)}}+|A|\big).
\]
Together with the bound on the differences this gives \eqref{eq:grid-domains}. If $|\nabla A|_{X^{(3)}}+[\nabla A]_{\alpha,X^{(3)}}<\delta'_1\le\delta'_0$ and $|A|<\delta'_0$, the right side of \eqref{eq:grid-domains} is less than $2\delta'_1+C_{\rm refl}\theta_r\delta'_0$. This proves the inclusion.
\end{proof}

\paragraph{The condition (DC6).} By (DC6) of Remark~\ref{rem:domain-constants}, $C_{\rm sym}\theta_r\delta_0\le\delta_1$, with $C_{\rm sym}=2C_{\rm refl}$ and $C_{\rm refl}$ of Lemma~\ref{lem:grid-domains}. Like (DC1)--(DC3), it is an upper bound on $\theta_r$, here in terms of $\delta_1/\delta_0$, with a constant that depends only on $d$. So it is compatible with them, and with the order of Remark~\ref{rem:parameters}(5), where $\theta_r$ is chosen after $\delta_1$ and $\delta_0$. The quantities that depend on the domain constants through $\theta_r/\delta_1$, such as $\Theta$ of Lemma~\ref{lem:later-normal-form}, only decrease. Under (DC6), Lemma~\ref{lem:grid-domains} with $\delta'_1=\delta_1^\flat$ and $\delta'_0=\delta_0$ gives $2\delta_1^\flat+C_{\rm refl}\theta_r\delta_0\le\frac12\delta_1+\frac12\delta_1$, so
\begin{equation}\label{eq:flat-in-full}
r\,\mathcal U^\flat(X)\subset\mathcal U(rX)\qquad\text{for every }r\in\mathfrak S^{\rm grid}_j\text{ and every polymer }X\text{ of scale }j.
\end{equation}
The term $e_{\mu\nu}$ in the proof of Lemma~\ref{lem:grid-domains} is of order $\eta\theta_r|A|$, and it does not vanish at curved real parts. Since $|A|$ is bounded only by $\delta_0$, it has to be absorbed in the budget $\delta_1$ of the differences; this is what (DC6) does. The factor $2$ in front of the H\"older quotient comes from Lemma~\ref{lem:collar-pairs}, and it is the reason for the factor $\frac14$ in $\delta_1^\flat$.

\begin{definition}[Raw generations and symmetrization]\label{def:symmetrization}
Let $h$ be an admissible threshold. A \emph{raw generation of scale $j$ at threshold $h$} is a family $E=\{E(X,\cdot)\}$, indexed by the polymers $X$ of scale $j$, with (G1) on the domains $\mathcal U(X)$ in place of $\mathcal U^\flat(X)$, and with (G3). That is, each $E(X,\cdot)$ is a normalized graded activity at $X$ on $\mathcal U(X)=\mathcal U(X;\theta_r,\delta_1,\delta_0)$, real at real configurations, and the function $\Phi_E$ of (G3) is invariant under $\mathfrak S_j$ on $\mathfrak U_j(h)$. Its norm is
\[
\|E\|^{\rm raw}_j=\sup_Xe^{\kappa_nd_j(X)}\sup_{\mathcal U(X)}|E(X,\cdot)| .
\]
Assume (DC6). The \emph{symmetrization} of $E$ is the family $E^{\rm sym}=\{E^{\rm sym}(X,\cdot)\}$ with
\[
E^{\rm sym}(X,U)=\frac1{|\mathfrak S^{\rm grid}_j|}\sum_{r\in\mathfrak S^{\rm grid}_j}E\big(r^{-1}X,r^{-1}U\big),\qquad U\in\mathcal U^\flat(X).
\]
\end{definition}

By \eqref{eq:flat-in-full}, applied to $r^{-1}$, $r^{-1}U\in\mathcal U(r^{-1}X)$ for $U\in\mathcal U^\flat(X)$, so every term is defined. A raw generation need not have (G2). The one-loop part $E'_{\rm ol}$ of a raw new generation, the sum of the pieces $\mathfrak e(X)$ and of the terms $\log\mathfrak j_c$ in Definition~\ref{def:raw-new} (Section~\ref{sec:ol-uv}), does have it, by Proposition~\ref{prop:one-loop-local}(b),(c) and Remark~\ref{rem:one-loop-local-use}(a); for the outputs $E^\#$ of the cluster expansion it is not proved. In the flow, the new activities of the step from scale $j-1$ to scale $j$, normalized and reblocked, form a raw generation (Lemma~\ref{lem:raw-new}, Theorem~\ref{thm:gen-step}(b)); (G3) for them is the symmetry statement of Theorem~\ref{thm:sf-step}, whose hypothesis is Remark~\ref{rem:later-symmetry}.

\begin{lemma}[Symmetrization]\label{lem:symmetrization}
Assume (DC6). Let $h$ be an admissible threshold, and let $E$ be a raw generation of scale $j$ at threshold $h$.
\begin{enumerate}[label=(\alph*)]
\item \emph{(Generation.)} $E^{\rm sym}$ is a generation of scale $j$ at threshold $h$: it has (G1)--(G4). Moreover $\|E^{\rm sym}\|_j\le\|E\|^{\rm raw}_j$.
\item \emph{(The sum at scale $j$.)} For every $V\in\mathfrak U_j(h)$,
\[
\sum_XE^{\rm sym}\big(X,U_j(V)\big)=\sum_XE\big(X,U_j(V)\big);
\]
that is, $\Phi_{E^{\rm sym}}=\Phi_E$.
\item \emph{(Later scales.)} Let $j\le k\le K$, and let $h''$ be an admissible threshold with $h\ge h''$ if $j=k$, and $h\ge2(C_*+1)L^{-2(k-j)}h''$ if $j<k$. Let $V\in\mathfrak U_k(h'')$ be a real configuration, in the units of scale $k$, and $U=U_k(V)$ an element of its minimal orbit. Then $U\in\mathcal U^\flat(X)$ for every polymer $X$ of scale $j$, in the units of scale $j$, and
\[
\sum_XE^{\rm sym}(X,U)=\sum_XE(X,U)=\Phi_E\big(U^{(j)}\big).
\]
Under the hypotheses of Lemma~\ref{lem:later-small-field}(c), $h=h_j$ and $h''=h_k$ are such thresholds.
\item \emph{(Complex points.)} Let $j\le k\le K$, let $\mathcal W$ be a finite-dimensional real vector space, let $\mathcal D$ be a connected open subset of its complexification, and let $B\mapsto U_B$ be a holomorphic map from $\mathcal D$ into the complex configurations on $\Tf$ such that $U_B\in\mathcal U^\flat(X)$ for every $B\in\mathcal D$ and every polymer $X$ of scale $j$. Suppose that $\mathcal D\cap\mathcal W$ contains a nonempty open subset $\mathcal D_\R$ of $\mathcal W$ such that, for $B\in\mathcal D_\R$, $U_B$ is an element of the minimal orbit of a real configuration $V_B$ as in (c). Then
\[
\sum_XE^{\rm sym}(X,U_B)=\sum_XE(X,U_B)\qquad\text{for every }B\in\mathcal D .
\]
\end{enumerate}
\end{lemma}

\begin{proof}\hasdetail{lem:symmetrization}
(a) \emph{(G1).} Fix $X$ and $r\in\mathfrak S^{\rm grid}_j$. Since $(r^{-1}U)(b)=U(rb)$, and $rb$ is a bond of $X$ or the reverse of one for every bond $b$ of $r^{-1}X$, the map $U|_X\mapsto(r^{-1}U)|_{r^{-1}X}$ is holomorphic. By \eqref{eq:flat-in-full}, it maps $\{U|_X:U\in\mathcal U^\flat(X)\}$ into $\{U'|_{r^{-1}X}:U'\in\mathcal U(r^{-1}X)\}$, on a neighbourhood of which $E(r^{-1}X,\cdot)$ is holomorphic. So every term of $E^{\rm sym}(X,\cdot)$ is a function of $U|_X$, holomorphic on a neighbourhood of $\{U|_X:U\in\mathcal U^\flat(X)\}$, and so is the finite average. For a real gauge transformation $u$, $r^{-1}(U^u)=(r^{-1}U)^{u\circ r}$, and $u\circ r$ is a real gauge transformation. Since $\mathcal U^\flat(X)$ is invariant under real gauge transformations and $E(r^{-1}X,\cdot)$ is invariant on $\mathcal U(r^{-1}X)$, $E^{\rm sym}(X,U^u)=E^{\rm sym}(X,U)$ for $U\in\mathcal U^\flat(X)$. Further $r^{-1}1=1$, so $E^{\rm sym}(X,1)=0$, and $r^{-1}U$ is real for real $U$. Finally $d_j(r^{-1}X)=d_j(X)$, so
\[
|E^{\rm sym}(X,U)|\le\max_r\sup_{\mathcal U(r^{-1}X)}|E(r^{-1}X,\cdot)|\le\|E\|^{\rm raw}_je^{-\kappa_nd_j(X)}\qquad(U\in\mathcal U^\flat(X)).
\]
This gives (G1) and the bound on the norm.

\emph{(G4) and (G2).} Let $s\in\mathfrak S^{\rm grid}_j$, and let $U\in\mathcal U^\flat(X)$ with $sU\in\mathcal U^\flat(sX)$. Substituting $r=st$,
\[
E^{\rm sym}(sX,sU)=\frac1{|\mathfrak S^{\rm grid}_j|}\sum_tE\big(t^{-1}s^{-1}sX,t^{-1}s^{-1}sU\big)=E^{\rm sym}(X,U).
\]
For a translation $s=\tau_a$, $a\in M\Z^4$, the condition $sU\in\mathcal U^\flat(sX)$ holds for all $U\in\mathcal U^\flat(X)$, and this is (G2). (G3) follows from (b).

(b) Let $V\in\mathfrak U_j(h)$, $U=U_j(V)$ and $r\in\mathfrak S^{\rm grid}_j$. Since $\mathfrak S^{\rm grid}_j\subset\mathfrak S_j$, Theorem~\ref{thm:balaban}(c) gives $r^{-1}V\in\mathfrak U_j(h)$, and $r^{-1}U$ is an element of the minimal orbit of $r^{-1}V$. The configurations $U$ and $r^{-1}U$ lie in every domain $\mathcal U^\flat(X)$ (the paragraph before Definition~\ref{def:generation}). The map $X\mapsto r^{-1}X$ is a bijection of the polymers of scale $j$. So, by (G3) for $E$,
\[
\sum_XE\big(r^{-1}X,r^{-1}U\big)=\sum_XE\big(X,r^{-1}U\big)=\Phi_E(r^{-1}V)=\Phi_E(V).
\]
Averaging over $r$ gives $\sum_XE^{\rm sym}(X,U)=\Phi_E(V)=\sum_XE(X,U)$. So $\Phi_{E^{\rm sym}}=\Phi_E$, which is invariant under $\mathfrak S_j$: this is (G3) for $E^{\rm sym}$.

(c) By Lemma~\ref{lem:later-small-field}(a) and the hypothesis on $h$, $U^{(j)}\in\mathfrak U_j(h)$, and by Lemma~\ref{lem:later-small-field}(b), $U$ is an element of the minimal orbit of $U^{(j)}$ for the problem of scale $j$. Since $h$ is admissible, the paragraph before Definition~\ref{def:generation} gives $U\in\mathcal U^\flat(X)\subset\mathcal U(X)$ for every $X$. By the definition of $\Phi_E$, with $U_j(U^{(j)})=U$, $\sum_XE(X,U)=\Phi_E(U^{(j)})$, and likewise for $E^{\rm sym}$. Now (b) gives the claim. The last sentence is Lemma~\ref{lem:later-small-field}(c).

(d) Each $E(X,\cdot)$ and each $E^{\rm sym}(X,\cdot)$ is holomorphic on a neighbourhood of $\{U|_X:U\in\mathcal U^\flat(X)\}$, the map $B\mapsto U_B|_X$ is holomorphic, and the sums are finite. So both sides are holomorphic on $\mathcal D$. By (c) they agree on $\mathcal D_\R$, a nonempty open subset of the totally real subspace $\mathcal W$ of maximal dimension. Since $\mathcal D$ is connected, the identity theorem gives the claim, as in the proof of Lemma~\ref{lem:complex-gauge}.
\end{proof}

\begin{remark}[The symmetrization in the flow]\label{rem:symmetrization-flow}
\begin{enumerate}[label=(\alph*)]
\item \emph{The density does not change.} The density of a step depends on a generation only through its sums at the background fields of small-field configurations (Remark~\ref{rem:flow-status}). By Lemma~\ref{lem:symmetrization}(b),(c), these are the sums of the raw generation. So the flow may carry $E^{\rm sym}$ in place of $E$ from the creation scale on; only the decomposition into activities changes. In a later step, Theorem~\ref{thm:sf-step} evaluates the old activities one at a time at the composed configurations $e^{\ii\eta A_*(s;\hat Z(Z))}U'$ of Proposition~\ref{prop:composed}. For $s\ne1$ these are not background fields, even for real $Z$, and no identity between the sums of $E$ and $E^{\rm sym}$ is claimed there. None is needed. The values of the old activities at $s\ne1$, at complex values of the class parameter $\lambda$ of Corollary~\ref{cor:two-classes}, and at complex references enter Sections~\ref{sec:step-quantities}--\ref{sec:cluster} only through bounds on each activity separately (Propositions~\ref{prop:composed} and~\ref{prop:localized-pieces}, Corollary~\ref{cor:two-classes}); for $E^{\rm sym}$ these bounds follow from $\|E^{\rm sym}\|_j\le\|E\|^{\rm raw}_j$. The identities that involve the density (Proposition~\ref{prop:small-field-step}, Theorem~\ref{thm:sf-step}, Remark~\ref{rem:extension}, the normalization at $U'=1$) are asserted only at real critical references and real $\xi$. The step is the integral of $e^{\delta\mathcal E(\xi)}$, and the localization reproduces $\delta\mathcal E(\xi)$ at $s=1$ (Section~\ref{subsec:cl-loc}). At a real critical reference and for real $\xi$, put $B=\hat Z(Z)$. The composed configuration at $s=1$ is $U_B=e^{\ii\eta A_*(B)}U'$ (Theorem~\ref{thm:weak-bg}(c)). Its real gauge transform $U_B^{w_B}$, $w_B=\mathsf w_k(e^{\ii\eta A_*(B)})$, is an element of the minimal orbit of $e^{\ii B}W$ (Theorem~\ref{thm:variational}(c), Corollary~\ref{cor:analytic}, Theorem~\ref{thm:balaban}(a)), at which Lemma~\ref{lem:symmetrization}(c) applies. The transformation $w_B$ need not lie in $\mathfrak G_0$. The domains $\mathcal U^\flat(X)$ and the activities of $E$ and $E^{\rm sym}$ are invariant under real gauge transformations, so their sums agree also at $U_B$. So the step gives the same density with $E$ and with $E^{\rm sym}$. The new activities it produces depend on which family is carried; the flow uses those produced from the symmetrized families.
\item \emph{Complex fields at $s=1$.} This is not needed for the flow, by (a); it shows that the sums of $E$ and $E^{\rm sym}$ agree also at complex points with $s=1$. Let $U'$ be a real critical reference of the step from scale $k\ge j$ to scale $k+1$, with base $W=U'^{(k)}$, and put $U_B=e^{\ii\eta A_*(1;B)}U'$ for coarse fields $B$ on the bonds of $\T^{(k)}$. For real $B$ with $|B|_\infty\le\min(w,R_P,\beta_V)$, $A_*(1;B)=A_*(B)$ (Theorem~\ref{thm:weak-bg}(c)), so $U_B^{w_B}$, with the real gauge transformation $w_B=\mathsf w_k(e^{\ii\eta A_*(B)})$ of $\Tf$, is an element of the critical orbit for $e^{\ii B}W$ (Corollary~\ref{cor:analytic}, Theorem~\ref{thm:variational}(c)). Moreover $e^{\ii B}W\in\mathfrak U_k(h_k)$ by (Q2), as in the proof of Proposition~\ref{prop:small-field-step}, so this orbit is its minimal orbit (Theorem~\ref{thm:balaban}(a)). The transformation $w_B$ need not lie in $\mathfrak G_0$. Put $\hat V_B=U_B^{(k)}$, the transform of $e^{\ii B}W$ by the restriction of $w_B^{-1}$ to $\T^{(k)}$. It is real, and $\hat V_B\in\mathfrak U_k(h_k)$, since the curvature is gauge invariant. By Theorem~\ref{thm:variational}(g) and Theorem~\ref{thm:balaban}(a), $U_B$ is an element of the minimal orbit of $\hat V_B$. For complex $B$ with $|B|_\infty<R_P$ and $C_{\rm conv}C_P|B|_\infty<\delta_1^\flat$, $U_B$ has real part $U'$ and a field bounded as in Definition~\ref{def:graded} with the budgets $\delta_1^\flat$ and $\delta_0$ (Theorem~\ref{thm:weak-bg}(a),(b) and Lemma~\ref{lem:norm-conversion}); so it lies in the domains $\mathcal U^\flat$ of scale $k$, and, by Remark~\ref{rem:domain-constants}(i), in the domains $\mathcal U^\flat(X)$ of scale $j$. On this polydisc, and under the hypotheses of Lemma~\ref{lem:later-small-field}(c) on the couplings, Lemma~\ref{lem:symmetrization}(d), with $V_B=\hat V_B$ for real $B$, shows that the sums of $E$ and $E^{\rm sym}$ agree at $U_B$, at complex $B$ as well.
\item \emph{The slack.} The newest generation, of scale $j$, is an old activity of the step from scale $j$ to scale $j+1$, on the domains $\mathcal U^\flat(X)$ of scale $j$. The reference is a point of the domain $\mathcal U$ of scale $j+1$, and on $X^{(3)}$ its representation has the budgets $L^{-2}\delta_1$ and $L^{-1}\delta_0$ (Remark~\ref{rem:domain-constants}(i)). So Proposition~\ref{prop:composed} applies with $\delta_1^\flat$ in place of $\delta_1$, and
\[
\min\big(\delta_1^\flat-L^{-2}\delta_1,\ \delta_0-L^{-1}\delta_0\big)\ge\big(\tfrac14-\tfrac19\big)\delta_1\ge\tfrac18\delta_1,
\]
since $L\ge3$ and $\delta_1\le\delta_0$. This replaces the slack $\frac12\delta_1$ of the domains $\mathcal U$ by $\frac18\delta_1$; it multiplies $\tau$ by $4$, and (Q3) remains a condition $\tau\le C_PL^dg_kp_{0,k}/\delta_1<\frac12$ with a larger $C_P$. At later scales $k>j$ the generation enters the step through the activities $\Psi_\Box$ and $E^{\rm ne}_Y$ of Section~\ref{sec:extraction}, which are bounded on the domains $\mathcal U^\flat$ of scale $k$; so the slack is again $\frac18\delta_1$. The single activities $E(X,\cdot)$ are defined on the whole domain $\mathcal U$ of scale $k$, since its points lie in $\mathcal U(X;L^{-2(k-j)}\theta_r,L^{-2(k-j)}\delta_1,L^{-(k-j)}\delta_0)\subset\mathcal U^\flat(X)$ (Remark~\ref{rem:domain-constants}(i)).
\item \emph{Why the domain shrinks.} Definition~\ref{def:graded} uses forward differences and coordinate-ordered paths, and a reflection turns forward differences into backward ones. At a curved real part a backward difference is a transported forward difference at a shifted point, up to the plaquette term $e_{\mu\nu}$ of Lemma~\ref{lem:grid-domains}; and the shifted points of a cube $\Lam_3$ of $(rX)^{(3)}$ need not lie in a cube $\Lam_3$ of $X^{(3)}$ near the outer boundary, which costs the factor $2$ of Lemma~\ref{lem:collar-pairs}. A symmetrized H\"older seminorm in Definition~\ref{def:graded} would remove the factor $2$, but not the plaquette term.
\end{enumerate}
\end{remark}

\subsection{The period density}\label{subsec:gen-period}

The coefficient of the counterterm of a generation is read off from its response to a constant curvature. For each polymer this response is a quadratic form $c^X$ in the curvature. Summed over a period of the grid, with the weights $1/|X|$, it gives the period density $\bar T_E$. By (G4) the period density is invariant under $B_4$, and the $B_4$-invariant forms are the multiples of the density of the action (Lemma~\ref{lem:b4-forms}). The comparison of the sum of a generation over the small polymers of a cube of a later scale with this multiple of the action of the cube is Section~\ref{sec:extraction}.

\paragraph{Curvature data and invariant forms.} A \emph{curvature datum} $F=(F_{\mu\nu})_{\mu<\nu}$ in $\mathfrak g^c$ is extended antisymmetrically, $F_{\nu\mu}=-F_{\mu\nu}$ and $F_{\mu\mu}=0$, and $|F|=\max_{\mu,\nu}|F_{\mu\nu}|$. For $R\in B_4$ put $R\cdot F=RFR^{\sf T}$, i.e.\ $(R\cdot F)_{\mu\nu}=\sum_{\kappa,\lambda}R_{\mu\kappa}R_{\nu\lambda}F_{\kappa\lambda}$, and for $v\in G$ let $vFv^{-1}$ be the curvature datum $(vF_{\mu\nu}v^{-1})$. With the bilinear form $\operatorname{tr}(XY)$ on $\mathfrak g^c$ (Section~\ref{sec:background}) put
\[
\operatorname{Iso}(F)=\sum_{\mu<\nu}\operatorname{tr}\big(F_{\mu\nu}^2\big),\qquad\operatorname{Pf}(F)=\operatorname{tr}\big(F_{12}F_{34}-F_{13}F_{24}+F_{14}F_{23}\big).
\]
A \emph{quadratic form} on the curvature data is a homogeneous polynomial of degree $2$ in their complex coordinates. Let $B_4^+=\{R\in B_4:\det R=1\}$.

\begin{lemma}[Invariant quadratic forms]\label{lem:b4-forms}
Let $G$ be compact, connected and simple, and let $\mathsf q$ be a quadratic form on the curvature data in $\mathfrak g^c$ with $\mathsf q(vFv^{-1})=\mathsf q(F)$ for all $v\in G$.
\begin{enumerate}[label=(\alph*)]
\item If $\mathsf q(R\cdot F)=\mathsf q(F)$ for all $R\in B_4$, then $\mathsf q=\beta\operatorname{Iso}$ for a unique $\beta\in\C$.
\item If $\mathsf q(R\cdot F)=\mathsf q(F)$ for all $R\in B_4^+$, then $\mathsf q=\beta\operatorname{Iso}+\tau\operatorname{Pf}$ for unique $\beta,\tau\in\C$. Both forms are invariant under $B_4^+$ and under conjugation, and $\operatorname{Pf}(R\cdot F)=(\det R)\operatorname{Pf}(F)$ for $R\in B_4$.
\item If $\mathsf q$ is real on the curvature data in $\mathfrak g$, then $\beta$ and $\tau$ are real.
\end{enumerate}
\end{lemma}

\begin{proof}\hasdetail{lem:b4-forms}
Let $\mathsf b$ be the symmetric bilinear form with $\mathsf b(F,F)=\mathsf q(F)$. Write $F=\sum_{\mu<\nu}e_{\mu\nu}\otimes F_{\mu\nu}$, with the basis $e_{\mu\nu}=e_\mu\wedge e_\nu$ of $\Lambda^2\C^4$; then $R\cdot F$ is the action of $\Lambda^2R$ on the first factor, $Re_{\mu\nu}=Re_\mu\wedge Re_\nu$. Since $G$ is connected, invariance under conjugation by $G$ is invariance under $\operatorname{ad}\mathfrak g$, hence under $\operatorname{ad}\mathfrak g^c$, as $\mathsf b$ is complex bilinear. So, for basis vectors $e,e'$ of $\Lambda^2\C^4$, $(X,Y)\mapsto\mathsf b(e\otimes X,e'\otimes Y)$ is an invariant bilinear form on $\mathfrak g^c$. Since $\mathfrak g$ is compact and simple, $\mathfrak g^c$ is simple, and its adjoint representation is irreducible. The form $\operatorname{tr}(XY)$ is invariant, and it is nondegenerate on $\mathfrak g^c$, since $\operatorname{tr}(X^2)=|X|^2>0$ for $0\ne X\in\mathfrak g$. Through it, an invariant bilinear form corresponds to an endomorphism of $\mathfrak g^c$ that commutes with $\operatorname{ad}\mathfrak g^c$, which is a multiple of the identity by Schur's lemma. Hence $\mathsf b(e\otimes X,e'\otimes Y)=Q(e,e')\operatorname{tr}(XY)$ with a bilinear form $Q$ on $\Lambda^2\C^4$, symmetric since $\mathsf b$ is, and
\[
\mathsf q(F)=\sum_{\mu<\nu}\ \sum_{\kappa<\lambda}Q(e_{\mu\nu},e_{\kappa\lambda})\operatorname{tr}\big(F_{\mu\nu}F_{\kappa\lambda}\big).
\]
So $\mathsf q$ is invariant under $R$ if and only if $Q(Rw,Rw')=Q(w,w')$ for all $w,w'$. The form Iso corresponds to the identity $I$, and Pf to the form $P$ with $P(e_{12},e_{34})=P(e_{14},e_{23})=\frac12$, $P(e_{13},e_{24})=-\frac12$, and $P=0$ on the other pairs of basis vectors.

\emph{Sign changes.} For $\epsilon\in\{\pm1\}^4$ the diagonal matrix $D_\epsilon$ maps $e_{\mu\nu}$ to $\epsilon_\mu\epsilon_\nu e_{\mu\nu}$, so invariance under $D_\epsilon$ gives $Q(e_{\mu\nu},e_{\kappa\lambda})=\epsilon_\mu\epsilon_\nu\epsilon_\kappa\epsilon_\lambda Q(e_{\mu\nu},e_{\kappa\lambda})$. Let $\{\mu,\nu\}\ne\{\kappa,\lambda\}$. If the two pairs share one index $\sigma$, let $\varrho$ be an index of one of them that is not in the other, and let $\epsilon_\sigma=\epsilon_\varrho=-1$, $\epsilon=1$ otherwise. Then $\det D_\epsilon=1$ and the product of the signs is $-1$; so $Q(e_{\mu\nu},e_{\kappa\lambda})=0$, already under $B_4^+$. If the pairs are disjoint, the product is $\det D_\epsilon$. Under $B_4$, a single sign change gives $Q(e_{\mu\nu},e_{\kappa\lambda})=0$. Under $B_4^+$ three entries remain: $p_1=Q(e_{12},e_{34})$, $p_2=Q(e_{13},e_{24})$ and $p_3=Q(e_{14},e_{23})$.

\emph{Permutations.} A permutation matrix maps $e_{\mu\nu}$ to $\pm e_{\pi(\mu)\pi(\nu)}$, with the indices in either order. Composed with a sign change if its determinant is $-1$, it lies in $B_4^+$ and acts on the pairs in the same way, up to signs. Since $S_4$ acts transitively on the pairs, the diagonal entries $Q(e_{\mu\nu},e_{\mu\nu})$ are all equal, to $\beta$ say. Under $B_4$ this gives $Q=\beta I$, which is (a). Under $B_4^+$, the even permutation matrix with $e_1\mapsto e_1$, $e_2\mapsto e_3$, $e_3\mapsto e_4$, $e_4\mapsto e_2$ maps the pair $(e_{12},e_{34})$ to $(e_{13},-e_{24})$ and $(e_{13},e_{24})$ to $(e_{14},-e_{23})$. So $p_1=-p_2$ and $p_2=-p_3$, and $Q=\beta I+2p_1P$, i.e.\ $\mathsf q=\beta\operatorname{Iso}+\tau\operatorname{Pf}$ with $\tau=2p_1$.

\emph{The two forms.} Both are invariant under conjugation. Iso is invariant under $B_4$, since every element of $B_4$ maps the basis $(e_{\mu\nu})$ onto itself up to signs. Under a sign change $D_\epsilon$, every term of Pf gets the factor $\epsilon_1\epsilon_2\epsilon_3\epsilon_4=\det D_\epsilon$. Under the transposition of $e_1$ and $e_2$, $(F_{12},F_{34},F_{13},F_{24},F_{14},F_{23})$ becomes $(-F_{12},F_{34},F_{23},F_{14},F_{24},F_{13})$, and Pf changes sign, since $\operatorname{tr}(XY)=\operatorname{tr}(YX)$; by symmetry the same holds for every transposition. Since $B_4$ is generated by the sign changes and the transpositions, $\operatorname{Pf}(R\cdot F)=(\det R)\operatorname{Pf}(F)$. For the curvature datum with the single entry $F_{12}=X\in\mathfrak g$, $X\ne0$, $\operatorname{Iso}(F)=\operatorname{tr}X^2>0$ and $\operatorname{Pf}(F)=0$; for $F_{12}=F_{34}=X$ and the other entries $0$, $\operatorname{Iso}(F)=2\operatorname{tr}X^2$ and $\operatorname{Pf}(F)=\operatorname{tr}X^2$. This gives the uniqueness in (a) and (b), and, evaluating $\mathsf q$ at these two real curvature data, (c).
\end{proof}

\paragraph{Lifts and linear fields.} Let $X$ be a polymer of scale $j$ that is not wrapping. As in the construction of the base cubes (Section~\ref{subsec:flow-norm}), the projection to a coordinate circle of the union of the closed cubes of the continuum filled by the cubes of $X$ is an arc. Let $\ell$ be its length. The points of the fine lattice in this arc are coordinates of points of $X$, and the two of them nearest to its ends are at distance $\ell-\eta$ along the arc. If $\ell-\eta>\frac12n_j$, two of these points are at distance at least $\frac12(n_j-\eta)$ on the circle; if $\ell-\eta\le\frac12n_j$, the two end points are at distance $\ell-\eta$ on the circle. Since $\operatorname{diam}X\le\frac12n_j-8$, the first case does not occur, and $\ell\le\frac12n_j-8+\eta$. The four arcs, enlarged by $3$ at both ends, have length less than $\frac12n_j$, and their product contains $X^{(3)}$. This product is the injective image of a box of $\R^4$ under the projection $\R^4\to\R^4/n_j\Z^4$, and we call the preimage of $X^{(3)}$ in this box a \emph{lift} of $X^{(3)}$. Lifts are unique up to translations by $n_j\Z^4$, the bonds of $X^{(3)}$ lift to bonds of $\eta\Z^4$, and the $\ell^\infty$-distance of two points of $X^{(3)}$ is that of their lifts. By Lemma~\ref{lem:torus}(b), whose first claim uses no condition on the torus, the points of $X$ are at mutual distance at most $M(2+d_j(X))$. So every point of a lift of $X^{(3)}$ is at $\ell^\infty$-distance at most
\[
D_X=M\big(2+d_j(X)\big)+3\le3M\big(1+d_j(X)\big)
\]
from the lift of any point of $X$.

For a curvature datum $F$ and a point $x_0\in\R^4$, the \emph{linear field} of $F$ with origin $x_0$ is the field on the oriented bonds of $\eta\Z^4$
\[
A_F(b)=\tfrac12\,(m_b-x_0)^{\sf T}F\,t_b=\tfrac12\sum_{\mu,\nu}(m_b-x_0)_\mu F_{\mu\nu}(t_b)_\nu ,
\]
where $m_b$ is the midpoint of $b$ and $t_b\in\{\pm e_1,\dots,\pm e_4\}$ its direction; so $A_F(-b)=-A_F(b)$. For $b=(b_-,b_-+\eta e_\nu)$, $A_F(b)=\frac12\sum_\mu F_{\mu\nu}(b_--x_0)_\mu$, since $F_{\nu\nu}=0$; for $x_0=x_X$ this is the field $A_{\rm lin}[F]$ of Section~\ref{sec:extraction}. Its differences are constant, $\partial_\mu A_{F,\nu}=\frac12F_{\mu\nu}$. Let $\tilde r(x)=Rx+\tilde a$ be an affine map of $\R^4$ with $R\in B_4$. The bond $\tilde r^{-1}b$ has midpoint $R^{-1}(m_b-\tilde rx_0)+x_0$ and direction $R^{-1}t_b$, and $R^{-1}=R^{\sf T}$; so the field $(\tilde rA_F)(b)=A_F(\tilde r^{-1}b)$ is
\begin{equation}\label{eq:linear-field-covariance}
\tilde rA_F=A_{R\cdot F}\qquad\text{with the origin }\tilde rx_0 .
\end{equation}
For a lift of $X^{(3)}$, an origin $x_0$ and $t\in\C$, let $V^X_{t,F}$ be the configuration equal to $e^{\ii\eta tA_F}$ on the bonds of $X^{(3)}$, read in the lift, and to $1$ on the other bonds.

\begin{definition}[Period density]\label{def:period-density}
Let $E=\{E(X,\cdot)\}$ be a family of normalized graded activities $E(X,\cdot)$ at the polymers $X$ of scale $j$ on the domains $\mathcal U^\flat(X)$, real at real configurations; for instance a generation of scale $j$, or a raw generation (Definition~\ref{def:symmetrization}). For a polymer $X$ of scale $j$ that is not wrapping and a curvature datum $F$ in $\mathfrak g^c$ put
\[
c^X(F)=\frac12\,\frac{d^2}{dt^2}\,E\big(X,V^X_{t,F}\big)\Big|_{t=0},
\]
for a lift of $X^{(3)}$ and an origin $x_0$. The \emph{period density} of $E$ is
\[
\bar T_E(F)=M^{-4}\sum_{X\supset Q}\frac{c^X(F)}{|X|},
\]
the sum running over the polymers $X\supset Q$ of scale $j$ that are not wrapping, where $Q$ is a grid cube of scale $j$ and $|X|=|X|_M$ is the number of grid cubes of $X$. The \emph{period density of the action} is
\[
\bar T_S(F)=n_j^{-4}\sum_pc^p(F),\qquad c^p(F)=\frac12\,\frac{d^2}{dt^2}\,s_p\big(V^p_{t,F}\big)\Big|_{t=0},
\]
the sum running over the plaquettes $p$ of $\Tf$, where $V^p_{t,F}$ equals $e^{\ii\eta tA_F}$ on the bonds of $p$, in coordinates of a neighbourhood of $p$, and $1$ elsewhere. If $\bar T_E=\beta\,\bar T_S$ for a number $\beta$, we write $\beta_E=\beta$.
\end{definition}

By Lemma~\ref{lem:period-density} below, $V^X_{t,F}\in\mathcal U^\flat(X)$ for small $|t|$, $c^X(F)$ depends only on $X$ and $F$, and $\bar T_E$ does not depend on $Q$ if $E$ has (G2). The weights $1/|X|$ make the reflections about the centre of $Q$ permute the terms of the sum (proof of Proposition~\ref{prop:period-iso}); the base cube of a polymer is not covariant under them. With (G2), the sum of $c^X$ over the polymers with base cube $Q$ is the same, $M^4\bar T_E$ (Section~\ref{sec:extraction}). The factor $n_j^{-4}$ is the inverse of the volume of the torus in the units of scale $j$, so $\bar T_S$ is the density per unit volume; by Lemma~\ref{lem:period-density}(c), $\bar T_E$ is such a density as well.

\begin{lemma}[The period density]\label{lem:period-density}
Let $G$ be semisimple, and let $E$ be as in Definition~\ref{def:period-density}.
\begin{enumerate}[label=(\alph*)]
\item Let $X$ be a polymer of scale $j$ that is not wrapping, and $F\ne0$ a curvature datum in $\mathfrak g^c$. Then $V^X_{t,F}\in\mathcal U^\flat(X)$ for small $|t|$, and for $|t|<\delta_1^\flat/(dD_X|F|)$ if the origin $x_0$ is the lift of a point of $X$. The number $c^X(F)$ does not depend on the lift of $X^{(3)}$ or on the origin $x_0$. It is a quadratic form in $F$, $c^X(vFv^{-1})=c^X(F)$ for $v\in G$, and
\[
|c^X(F)|\le9d^2M^2\big(1+d_j(X)\big)^2(\delta_1^\flat)^{-2}\,\|E(X)\|\,|F|^2,\qquad\|E(X)\|=\sup_{\mathcal U^\flat(X)}|E(X,\cdot)|.
\]
\item If $E$ has (G2), then $c^{X+a}=c^X$ for $a\in M\Z^4$. If $E$ has (G4), then $c^{rX}(R\cdot F)=c^X(F)$ for every $r\in\mathfrak S^{\rm grid}_j$, with $R$ the linear part of $r$.
\item Let $\|E\|_j=\sup_Xe^{\kappa_nd_j(X)}\|E(X)\|<\infty$. The sum defining $\bar T_E(F)$ converges absolutely, uniformly in the torus, and
\[
|\bar T_E(F)|\le9d^2K_0M^{-2}(\delta_1^\flat)^{-2}\,\|E\|_j\,|F|^2 ,
\]
with $K_0$ of \eqref{eq:dM-facts}. If $E$ has (G2), then $\bar T_E$ does not depend on $Q$, and
\[
\bar T_E(F)=n_j^{-4}\sum_Xc^X(F),
\]
the sum running over all polymers of scale $j$ that are not wrapping.
\item For $p=p_{\mu\nu}(x)$, $\mu<\nu$, $c^p(F)=\frac12\eta^4\operatorname{tr}(F_{\mu\nu}^2)$. Hence
\[
\bar T_S(F)=\tfrac12\sum_{\mu<\nu}\operatorname{tr}\big(F_{\mu\nu}^2\big)=\tfrac12\operatorname{Iso}(F).
\]
Let $\mathcal P(X)$, $X$ a polymer of scale $j$, be sets of plaquettes such that every plaquette lies in exactly one of them, $\mathcal P(X)=\emptyset$ for wrapping $X$, and the bonds of the plaquettes of $\mathcal P(X)$ are bonds of $X$. If the family $E(X,\cdot)=\sum_{p\in\mathcal P(X)}s_p$ is as in Definition~\ref{def:period-density} and has (G2), then $\bar T_E=\bar T_S$.
\end{enumerate}
\end{lemma}

\begin{proof}\hasdetail{lem:period-density}
(a) For a field $B$ on the bonds of a lift of $X^{(3)}$, let $V_B$ equal $e^{\ii\eta B}$ on the bonds of $X^{(3)}$, read in the lift, and $1$ elsewhere. With the real part $1$ and the field $B$ in Definition~\ref{def:graded}, $V_B\in\mathcal U^\flat(X)$ if $|B|<\delta_0$ and $|\partial B|+[\partial B]_\alpha<\delta_1^\flat$ on $X^{(3)}$. This holds for $B=tA_F$ with $|t|$ small, since $A_F$ is bounded on the finite lift. If $x_0$ is the lift of a point of $X$ and $|t|<\delta_1^\flat/(dD_X|F|)$, then $|tA_F|\le\frac12(d-1)D_X|F||t|<\delta_1^\flat\le\delta_0$ and $|\partial(tA_F)|=\frac12|t||F|<\delta_1^\flat$, and the H\"older quotients vanish. This gives the first claim. So $\Phi_X(B)=E(X,V_B)$ is holomorphic for small $B$, and $c^X(F)=\frac12d^2\Phi_X(0)[A_F,A_F]$ is a quadratic form in $F$, since $A_F$ is linear in $F$.

As in Step~3 of the proof of Lemma~\ref{lem:quadratic-bound}, $d\Phi_X(0)=0$, since $G$ is semisimple: by real gauge invariance under constant gauge transformations $v\in G$, the gradient of $\Phi_X$ at $0$ is $\operatorname{Ad}(G)$-invariant at every bond, and for semisimple $G$ the only such vector is $0$. Let $\lambda$ be a small real scalar field on $\Tf$ and $u=e^{\ii\lambda}$. On the bonds of $X^{(3)}$, $V_B^u=V_{B^\lambda}$, where $e^{\ii\eta B^\lambda(b)}=e^{\ii\lambda(b_-)}e^{\ii\eta B(b)}e^{-\ii\lambda(b_+)}$, and $B^\lambda=B-\partial\lambda+O(|B||\lambda|+|\lambda|^2)$ is analytic in $(B,\lambda)$, with $\partial\lambda(b)=\eta^{-1}(\lambda(b_+)-\lambda(b_-))$. Since $E(X,\cdot)$ depends only on the bonds of $X$ and is invariant under real gauge transformations on $\mathcal U^\flat(X)$, $\Phi_X(B^\lambda)=\Phi_X(B)$ for small $B$ and $\lambda$. Since $d\Phi_X(0)=0$, the terms of second order in $(B,\lambda)$ give $d^2\Phi_X(0)[B-\partial\lambda,B-\partial\lambda]=d^2\Phi_X(0)[B,B]$. Replacing $B$ by $sB$ and comparing the terms linear in $s$, $d^2\Phi_X(0)[B,\partial\lambda]=0$ for all $B$ and all real $\lambda$, hence, by linearity, for complex $\lambda$. Changing the origin from $x_0$ to $x_0+v$ adds to $A_F$ the field $-\frac12v^{\sf T}Ft_b=\partial\lambda$ with $\lambda(x)=-\frac12v^{\sf T}Fx$ on the lift. Hence
\[
\tfrac12d^2\Phi_X(0)[A_F+\partial\lambda,A_F+\partial\lambda]=\tfrac12d^2\Phi_X(0)[A_F,A_F],
\]
and $c^X(F)$ does not depend on the origin. Two lifts differ by a translation by a vector of $n_j\Z^4$, which amounts to a change of the origin. For constant $v\in G$, $V^X_{t,vFv^{-1}}=(V^X_{t,F})^v$, since $A_{vFv^{-1}}=vA_Fv^{-1}$ and the configuration $1$ is invariant; so real gauge invariance gives $c^X(vFv^{-1})=c^X(F)$. For the bound, take $x_0$ the lift of a point of $X$. Then $f(t)=E(X,V^X_{t,F})$ is holomorphic for $|t|<t_X=\delta_1^\flat/(dD_X|F|)$ and bounded there by $\|E(X)\|$, and $c^X(F)=\frac12f''(0)$. Cauchy's estimate gives $|c^X(F)|\le\|E(X)\|t_X^{-2}=d^2D_X^2(\delta_1^\flat)^{-2}\|E(X)\||F|^2$, and $D_X\le3M(1+d_j(X))$.

(b) Let $a\in M\Z^4$. Translating the lift and the origin by $a$, $\tau_aV^X_{t,F}=V^{X+a}_{t,F}$, and both configurations lie in the domains $\mathcal U^\flat$ for small $|t|$, by (a). So (G2) gives $E(X+a,V^{X+a}_{t,F})=E(X,V^X_{t,F})$, and $c^{X+a}=c^X$. Let $r\in\mathfrak S^{\rm grid}_j$, and let $\tilde r$ be the affine map of $\R^4$ that lifts $r$ and maps the given lift of $X^{(3)}$ onto a lift of $(rX)^{(3)}=r(X^{(3)})$. The configuration $rV^X_{t,F}$ has real part $1$, and on reversed bonds its field is $-A_F$; so \eqref{eq:linear-field-covariance} gives $rV^X_{t,F}=V^{rX}_{t,R\cdot F}$, for this lift and the origin $\tilde rx_0$. Since $|R\cdot F|=|F|$, both configurations lie in the domains $\mathcal U^\flat$ for small $|t|$, by (a). So (G4) gives $c^{rX}(R\cdot F)=c^X(F)$.

(c) By the bound of (a), with $\|E(X)\|\le\|E\|_je^{-\kappa_nd_j(X)}$ and $|X|\ge1$,
\[
\sum_{X\supset Q}\frac{|c^X(F)|}{|X|}\le9d^2M^2(\delta_1^\flat)^{-2}\|E\|_j|F|^2\sum_{X\supset Q}\big(1+d_j(X)\big)^2e^{-\kappa_nd_j(X)} .
\]
By Remark~\ref{rem:parameters}(1), $\kappa_n-\kappa_0\ge\kappa_0+2\ge3$, and $(1+t)^2e^{-3t}\le1$ for $t\ge0$. So the last sum is at most $\sum_{X\supset Q}e^{-\kappa_0d_j(X)}\le K_0$, by \eqref{eq:dM-facts}, uniformly in the torus. This gives the bound. If $E$ has (G2), the translation by $a\in M\Z^4$ maps the polymers $X\supset Q$ that are not wrapping bijectively onto those that contain $Q+a$, and preserves $|X|$; with (b), the sum for $Q+a$ equals that for $Q$. Since the translations by $M\Z^4$ act transitively on the grid cubes, $\bar T_E$ does not depend on $Q$. There are $(n_j/M)^4$ grid cubes, and every polymer has $|X|$ of them; so
\[
\sum_Xc^X(F)=\sum_X\ \sum_{Q\subset X}\frac{c^X(F)}{|X|}=\sum_Q\ \sum_{X\supset Q}\frac{c^X(F)}{|X|}=\Big(\frac{n_j}M\Big)^4M^4\,\bar T_E(F),
\]
all sums running over polymers that are not wrapping.

(d) Let $p=p_{\mu\nu}(x)$, with the bonds $b_1,\dots,b_4$ of Definition~\ref{def:curvature} on the fine lattice, and $a_l=A_F(b_l)$. Since the real part is $1$, $V^p_{t,F}(\partial p)=e^{\ii\eta ta_1}e^{\ii\eta ta_2}e^{-\ii\eta ta_3}e^{-\ii\eta ta_4}$. Here $a_2-a_4=\eta\partial_\mu A_{F,\nu}(x)=\frac12\eta F_{\mu\nu}$ and $a_1-a_3=-\eta\partial_\nu A_{F,\mu}(x)=\frac12\eta F_{\mu\nu}$, so $a_1+a_2-a_3-a_4=\eta F_{\mu\nu}$, whatever the origin. By the Baker--Campbell--Hausdorff formula, $V^p_{t,F}(\partial p)=e^{\ii Y(t)}$ with $Y(t)=t\eta^2F_{\mu\nu}+O(t^2)$. Since $s_p(U)=\mathsf N-\frac12\operatorname{tr}(U(\partial p)+U(\partial p)^{-1})$ (Section~\ref{sec:action}), $s_p(V^p_{t,F})=\operatorname{tr}(1-\cos Y(t))=\frac12\operatorname{tr}(Y(t)^2)+O(t^4)=\frac12t^2\eta^4\operatorname{tr}(F_{\mu\nu}^2)+O(t^3)$. So $c^p(F)=\frac12\eta^4\operatorname{tr}(F_{\mu\nu}^2)$. There are $(n_j/\eta)^4$ plaquettes in each plane, which gives $\bar T_S$. For the last claim, $E(X,V^X_{t,F})=\sum_{p\in\mathcal P(X)}s_p(V^X_{t,F})$, and on the bonds of $p\in\mathcal P(X)$ the configuration $V^X_{t,F}$ is $e^{\ii\eta tA_F}$ in coordinates of a neighbourhood of $p$; so $c^X_E=\sum_{p\in\mathcal P(X)}c^p$. By (c), $\bar T_E=n_j^{-4}\sum_X\sum_{p\in\mathcal P(X)}c^p=n_j^{-4}\sum_pc^p=\bar T_S$.
\end{proof}

\begin{proposition}[Isotropy of the period density]\label{prop:period-iso}
Let $G$ be compact, connected and simple, and let $E$ be a generation of scale $j$ (Definition~\ref{def:generation}); property (G3) is not used.
\begin{enumerate}[label=(\alph*)]
\item There is a unique number $\beta_E$ with $\bar T_E=\beta_E\bar T_S$, that is,
\[
\bar T_E(F)=\tfrac12\beta_E\sum_{\mu<\nu}\operatorname{tr}\big(F_{\mu\nu}^2\big)\qquad\text{for all curvature data }F\text{ in }\mathfrak g^c .
\]
It is real. So the period density of $E$ has no topological part, a multiple of Pf, and no anisotropic part.
\item $|\beta_E|\le C_\beta\|E\|_j$, with $C_\beta=18d^2K_0M^{-2}(\delta_1^\flat)^{-2}$.
\item Let $E=E'^{\rm sym}$ for a raw generation $E'$ (Definition~\ref{def:symmetrization}), and put $\bar T^{\rm vol}_{E'}(F)=n_j^{-4}\sum_Xc^X_{E'}(F)$, the sum running over the polymers that are not wrapping. Then
\[
\bar T_E(F)=\frac1{|B_4|}\sum_{R\in B_4}\bar T^{\rm vol}_{E'}(R\cdot F),\qquad\beta_E=\frac1{3\dim\mathfrak g}\sum_{\mu<\nu}\ \sum_a\bar T^{\rm vol}_{E'}\big(F^{\mu\nu,a}\big),
\]
where $(T_a)$ is a basis of $\mathfrak g$ with $\operatorname{tr}(T_aT_b)=\delta_{ab}$, and $F^{\mu\nu,a}$ is the curvature datum with the single entry $F_{\mu\nu}=T_a$.
\end{enumerate}
\end{proposition}

\begin{proof}\hasdetail{prop:period-iso}
(a) By (G2) and Lemma~\ref{lem:period-density}(c), $\bar T_E(F)=n_j^{-4}\sum_Xc^X(F)$, a quadratic form in $F$, invariant under conjugation by Lemma~\ref{lem:period-density}(a). Let $R\in B_4$, and let $r(x)=Rx$, an element of $\mathfrak S^{\rm grid}_j$. It maps the polymers that are not wrapping bijectively onto themselves, and $c^{rX}(R\cdot F)=c^X(F)$ by (G4) and Lemma~\ref{lem:period-density}(b). So
\[
\bar T_E(R\cdot F)=n_j^{-4}\sum_Xc^{rX}(R\cdot F)=n_j^{-4}\sum_Xc^X(F)=\bar T_E(F).
\]
The same follows from the definition with the weights $1/|X|$, for the grid cube $Q$ centred at $0$: $r$ maps the polymers $X\supset Q$ onto themselves, since $rQ=Q$, and preserves $|X|$. Lemma~\ref{lem:b4-forms}(a) gives $\bar T_E=\beta'\operatorname{Iso}$ for a unique $\beta'$, and $\bar T_S=\frac12\operatorname{Iso}$ by Lemma~\ref{lem:period-density}(d); so $\beta_E=2\beta'$. By (G1), $t\mapsto E(X,V^X_{t,F})$ is real for real $t$ and real curvature data $F$, so $c^X$ and $\bar T_E$ are real on real curvature data, and $\beta_E$ is real by Lemma~\ref{lem:b4-forms}(c).

(b) Let $F^*=F^{12,a}$ for some $a$. Then $|F^*|=|T_a|=1$, since $|X|^2=\operatorname{tr}(X^2)$ on $\mathfrak g$, and $\bar T_S(F^*)=\frac12$. So $\beta_E=2\bar T_E(F^*)$, and Lemma~\ref{lem:period-density}(c) gives the bound.

(c) Let $X$ be a polymer that is not wrapping, $r\in\mathfrak S^{\rm grid}_j$ with linear part $R_r$, and $\tilde r$ as in the proof of Lemma~\ref{lem:period-density}(b). As there, $r^{-1}V^X_{t,F}=V^{r^{-1}X}_{t,R_r^{-1}\cdot F}$, for the lift $\tilde r^{-1}$ of the given one and the origin $\tilde r^{-1}x_0$. So, for small $|t|$,
\[
E\big(X,V^X_{t,F}\big)=\frac1{|\mathfrak S^{\rm grid}_j|}\sum_rE'\big(r^{-1}X,V^{r^{-1}X}_{t,R_r^{-1}\cdot F}\big),\qquad c^X_E(F)=\frac1{|\mathfrak S^{\rm grid}_j|}\sum_rc^{r^{-1}X}_{E'}\big(R_r^{-1}\cdot F\big).
\]
Summing over $X$, and using that $X\mapsto r^{-1}X$ permutes the polymers that are not wrapping, $\sum_Xc^X_E(F)=|\mathfrak S^{\rm grid}_j|^{-1}\sum_r\sum_Xc^X_{E'}(R_r^{-1}\cdot F)$. The linear part $R_r$ runs over $B_4$, each element being taken by $(n_j/M)^4$ elements $r$, and $R\mapsto R^{-1}$ is a bijection of $B_4$. With Lemma~\ref{lem:period-density}(c) for $E$, which has (G2), this is the first formula. Every $R\in B_4$ maps each curvature datum $F^{\mu\nu,a}$ to $\pm F^{\mu'\nu',a}$, permuting the pairs $\{\mu,\nu\}$, and quadratic forms are even. So the first formula gives $\sum_{\mu<\nu,a}\bar T_E(F^{\mu\nu,a})=\sum_{\mu<\nu,a}\bar T^{\rm vol}_{E'}(F^{\mu\nu,a})$. By (a) the left side is $\frac12\beta_E\sum_{\mu<\nu,a}\operatorname{tr}(T_a^2)=3\dim\mathfrak g\,\beta_E$.
\end{proof}

\begin{remark}[On the period density]\label{rem:period-density}
\begin{enumerate}[label=(\alph*)]
\item \emph{The reflections are needed.} Under $B_4^+$ alone, Lemma~\ref{lem:b4-forms}(b) leaves a multiple $\tau\operatorname{Pf}$, the density of the topological charge. The symmetry (G3) of the sum of a family does not exclude a multiple of Pf in its period density, and the note does not show that its coefficient is controlled by the sum. A surviving $\tau$ would contribute terms of order $|\tau|$ per cube of every later scale, with no decay in the number of scales. The covariance (G4) under all of $\mathfrak S^{\rm grid}_j$, which includes the reflections, excludes it by Proposition~\ref{prop:period-iso}; this is why the generations are symmetrized over $B_4$, and not only over $B_4^+$ or the translations.
\item \emph{Scales.} The number $\beta_E$ does not depend on the units: in the units of another scale both $\bar T_E$ and $\bar T_S$ change by the same factor. It does not depend on the threshold of the generation.
\item \emph{Relation to Part~XI and to Ba\l aban's coefficient.} The relation of $\beta_E$, for the generations of the flow, to the one-loop coefficient of Part~XI and to the coefficient that Ba\l aban's definition \cite[(1.20)--(1.22)]{BalYM4I} assigns to the generation is (F2) of Proposition~\ref{prop:flow-F} and Proposition~\ref{prop:beta-identification}, the latter through the generation locality of Section~\ref{sec:gen-locality}, without a rate; neither is used in Part~X. These statements take into account that $\beta_E$ is a property of the family $E$, not shown to be determined by its sum $\Phi_E$ alone (Section~\ref{subsec:ol-period-vacpol}); for the flow this does not matter, since Proposition~\ref{prop:extraction} applies to every generation with (G1)--(G4). Proposition~\ref{prop:period-iso}(c) expresses $\beta_E$ for a symmetrized generation through the raw family produced by the step: it is the isotropic part of the volume average of the quadratic forms $c^X_{E'}$, and the Pfaffian and anisotropic parts of that average are removed by the symmetrization.
\end{enumerate}
\end{remark}
